%% DeReFusion 复现报告 —— 按 Applied Soft Computing (elsarticle) 格式规范生成
%% 编译：xelatex/pdflatex + elsarticle.cls（本目录附官方 cls 与 .bst）
%% 期刊格式要素：frontmatter（title/author/affiliation/abstract/keyword）+
%% graphicalabstract + highlights + CRediT + declaration + appendix + elsarticle-num 参考文献
\documentclass[preprint,12pt]{elsarticle}

\usepackage{amssymb}
\usepackage{amsmath}
\usepackage{booktabs}
\usepackage{graphicx}
\usepackage{float}
\usepackage{hyperref}

\journal{Applied Soft Computing}

\begin{document}

\begin{frontmatter}

\title{Reproduction Study of DeReFusion: A Controlled Comparison of Soft
Computing Fusion Strategies for Financial Time Series Forecasting\tnoteref{t1}}
\tnotetext[t1]{Reproduction of Hsieh and Chen (2026), Applied Soft Computing 203, 116252.}

\author[nwpu]{Dongxu Jiang\corref{cor1}}
\ead{jiangdongxu04@gmail.com}
\cortext[cor1]{Corresponding author. Email: jiangdongxu04@gmail.com.}
\affiliation[nwpu]{organization={},
            city={},
            country={}}

\begin{abstract}
This report documents a partial reproduction of the DeReFusion framework
proposed by Hsieh and Chen (2026), which compares fusion strategies for
financial time series forecasting under a controlled protocol: a DLinear base
branch, an LSTM--Transformer residual branch, and parameter-free additive
fusion, all wrapped in Reversible Instance Normalization (RevIN). Using the
authors' official code repository with newly fetched Yahoo Finance daily OHLC
data (2016--2025), we reproduce six controlled experiments on the S\&P 500
index: the full DeReFusion model, a RevIN-wrapped DLinear baseline, and a
learnable-gate fusion variant, at forecast horizons $T\in\{1,24\}$. The
central qualitative finding of the original paper is confirmed: under an
identical protocol, parameter-free additive fusion outperforms learnable
gating (MSE 0.0623 vs.\ 0.0719 at $T=24$), and the hybrid model is
competitive with the strongest linear baseline. The model size (28{,}436
trainable parameters) matches the paper's reported ${\sim}26$k. We further
distill the experimental protocol into six verified insertion points
(checkpoints) that allow new research directions---in particular,
structure-aware expert routing---to be evaluated within the same framework
at minimal engineering cost. All materials (forked code, dataset, fetch
scripts, and per-experiment artifacts) are available in a public
repository.
\end{abstract}

\begin{graphicalabstract}
% 论文要求：单独图形摘要文件；此处嵌入复现主图作为示意
\includegraphics[width=\textwidth]{figs/fig_derefusion_t24_dashboard.png}
\end{graphicalabstract}

\begin{highlights}
\item Full pipeline of DeReFusion reproduced from official code on new data
\item Additive fusion beats learnable gating under identical protocol
\item Model size 28k parameters matches the paper's reported scale
\item Six verified insertion points documented for follow-up experiments
\item Environment pitfalls (encoding, CPU mode, dependency pins) recorded
\end{highlights}

\begin{keyword}
Financial time series forecasting \sep Reproduction study \sep
Decomposition--residual architecture \sep Fusion strategy comparison \sep
Reversible instance normalization
\end{keyword}

\end{frontmatter}

% \linenumbers

\section{Introduction}
\label{sec:intro}

DeReFusion~\cite{hsieh2026derefusion} poses a deliberately simple question:
on non-stationary financial series, how far can a linear decomposition base,
a learned residual, and plain additive fusion go, and does anything more
elaborate buy accuracy? The paper answers it with a controlled benchmark: a
single data pipeline, normalization scheme, training protocol, and metric
set shared by every model, so that differences in results reflect the
backbone rather than the harness. The authors conclude that parameter-free
additive fusion beats static weighting and dynamic gating, that
architectural complexity is no guarantee of accuracy, and that statistical
accuracy does not translate into economic value.

This document serves two purposes. First, it reports a partial reproduction
of the paper's experiments from the official repository, verifying that the
published protocol and qualitative conclusions are attainable from the
public code with freshly fetched data. Second, it distills the experimental
machinery into a set of verified \emph{insertion points} so that follow-up
research---ours is structure-aware expert routing for financial panel
forecasting---can be evaluated inside the same framework with minimal
engineering.

The remainder of this report is organized as follows.
Section~\ref{sec:method} summarizes the reproduced architecture and
protocol. Section~\ref{sec:setup} describes the reproduction environment
and deviations from the original setup. Section~\ref{sec:results} presents
the reproduced results and compares them against the paper's reported
values. Section~\ref{sec:checkpoints} formalizes six verified insertion
points for future work, and Section~\ref{sec:lessons} records engineering
lessons. Section~\ref{sec:conclusion} concludes.

\section{Reproduced framework}
\label{sec:method}

\subsection{Architecture}

DeReFusion is a dual-branch decomposition--residual forecaster wrapped in
RevIN~\cite{kim2022revin}. Given an input window
$\mathbf{X}\in\mathbb{R}^{B\times L\times C}$ ($C=4$ OHLC channels):

\begin{enumerate}
  \item \textbf{RevIN normalization} strips instance-level mean/variance
        shift; statistics are retained for symmetric de-normalization.
  \item \textbf{Linear base branch (DLinear):} moving-average decomposition
        (kernel $k=25$) splits $\mathbf{X}$ into seasonal and trend
        components; channel-shared linear maps
        $W_{\mathrm{se}}, W_{\mathrm{tr}}\in\mathbb{R}^{T\times L}$ produce
        the base forecast $\mathbf{Y}_{\mathrm{base}}$.
  \item \textbf{Nonlinear residual branch:} per-feature projection
        $\mathbb{R}^{C\to d}$, a one-layer LSTM, a one-layer Transformer
        encoder (4 heads, FFN width $4d$, dropout 0.3), then temporal
        ($\mathbb{R}^{L\to T}$) and channel ($\mathbb{R}^{d\to C}$)
        projections yield $\mathbf{Y}_{\mathrm{res}}$.
  \item \textbf{Direct additive fusion:}
        $\mathbf{F}=\mathbf{Y}_{\mathrm{base}}+\mathbf{Y}_{\mathrm{res}}$
        with no fusion parameters.
  \item \textbf{RevIN de-normalization} restores the original scale.
\end{enumerate}

\subsection{Controlled protocol}

All models share: chronological 70/10/20 splitting; StandardScaler fitted on
the training split only; RevIN on every trained model; multivariate input
with single-target output (MS; target \texttt{Close}); Adam with learning
rate $10^{-4}$, batch size 32, MSE loss, early stopping (patience 5) on
validation MSE, cosine schedule; random seed 2021. Metrics are MAE, MSE,
RMSE, MAPE, MSPE, and $R^2$, computed on the normalized series.

\section{Reproduction setup and deviations}
\label{sec:setup}

Table~\ref{tab:setup} lists the setup and deviations from the original
study. Deviations are limited to factors that do not affect the protocol's
fairness (single asset, single seed, CPU training, newer library versions).

\begin{table}[H]
\centering
\caption{Reproduction setup and deviations from the original study.}
\label{tab:setup}
\begin{tabular}{lll}
\toprule
Item & Original & This reproduction \\
\midrule
Code & Official repo (MIT) & Fork of official repo \\
Data & 10 assets, 2016--2025, Yahoo & GSPC only, Yahoo, 2016--2025 (2{,}513 rows) \\
Assets grid & 10 assets $\times$ 5 horizons & 1 asset $\times$ 2 horizons \\
Horizons $T$ & $\{1,7,12,24,36\}$ & $\{1,24\}$ \\
Seeds & 5 seeds (2020--2024) & 1 seed (2021) \\
Models & 16 (10 RevIN baselines + gates + z.s.) & 3 (DeReFusion, revin-DLinear, gatev2) \\
Hardware & 3 machines, CUDA & 1 laptop, CPU-only \\
PyTorch & 2.5.1 & 2.14.0 (CPU) \\
Python & 3.11 & 3.13.5 \\
\bottomrule
\end{tabular}
\end{table}

The dataset was regenerated from Yahoo Finance with a purpose-written fetch
script (\texttt{scripts/fetch\_dataset\_v3.py}), since the upstream
repository ignores \texttt{dataset/} in version control. Row counts match
the paper for every asset class (2{,}513 trading days for equities/indices;
2{,}602 for FX; 3{,}653/2{,}975 for BTC/ETH).

\section{Reproduced results}
\label{sec:results}

\subsection{Main comparison}

Table~\ref{tab:main} reports the reproduced test metrics on GSPC with seed
2021. Values are MSE/R$^2$ on the normalized series, consistent with the
paper's convention.

\begin{table}[H]
\centering
\caption{Reproduced test metrics (GSPC, seed 2021). Best value per column in bold.}
\label{tab:main}
\begin{tabular}{lcccc}
\toprule
Model & MSE ($T{=}1$) & MSE ($T{=}24$) & $R^2$ ($T{=}24$) & Params \\
\midrule
DeReFusion (additive)   & 0.0165 & \textbf{0.0623} & \textbf{0.869} & 28{,}436 \\
revin-DLinear           & \textbf{0.0158} & 0.0701 & 0.853 & ${\sim}3$k \\
gatev2-learnable        & 0.0253 & 0.0719 & 0.849 & ${\sim}28$k \\
\midrule
\multicolumn{5}{l}{\footnotesize Paper anchors (10-asset means): DeReFusion $T{=}1$ $\approx$ 0.037, $T{=}24$ $\approx$ 0.205; DLinear $\approx$ 0.049/0.172.}\\
\end{tabular}
\end{table}

\subsection{Verification against the paper's conclusions}

\paragraph{Additive fusion beats learnable gating.} At $T=24$ the ordering
DeReFusion $(0.0623) <$ DLinear $(0.0701) <$ gatev2 $(0.0719)$ matches the
paper's central claim that parameter-free additive fusion outperforms
parameterized gating. The gate variant is worse than even the linear
baseline, echoing the paper's finding that gates under-fit on
sample-limited, low signal-to-noise financial series.

\paragraph{Model scale.} The reproduced DeReFusion has 28{,}436 trainable
parameters, in line with the paper's ${\sim}26$k (difference attributable
to the horizon-dependent output head).

\paragraph{Hybrid vs.\ linear at short horizons.} On this single asset the
DLinear baseline marginally outperforms DeReFusion at $T=1$ (0.0158 vs.\
0.0165, a 4\% gap). The paper, averaging over ten assets, reports a 24\%
advantage for DeReFusion at $T=1$ while also stressing that differences
inside the leading cluster are tiny and asset-dependent. Our single-asset,
single-seed observation therefore does not contradict the paper; it
illustrates exactly the variance structure the paper acknowledges.

\subsection{Qualitative artifacts}

Figures~\ref{fig:deref} and \ref{fig:curves} show the reproduced dashboard
(six-metric radar, error analysis, prediction curves) generated by the
repository's own visualization utilities from saved predictions.

\begin{figure}[H]
\centering
\includegraphics[width=0.85\textwidth]{figs/fig_derefusion_t24_dashboard.png}
\caption{DeReFusion, GSPC, $T=24$: reproduction dashboard (metrics radar,
error analysis, error heatmap).}
\label{fig:deref}
\end{figure}

\begin{figure}[H]
\centering
\includegraphics[width=0.85\textwidth]{figs/fig_derefusion_t24_curves.png}
\caption{DeReFusion, $T=24$: sample ground-truth vs.\ prediction windows.}
\label{fig:curves}
\end{figure}

\begin{figure}[H]
\centering
\includegraphics[width=0.48\textwidth]{figs/fig_dlinear_t24_dashboard.png}\hfill
\includegraphics[width=0.48\textwidth]{figs/fig_gatev2_t24_dashboard.png}
\caption{Baseline dashboards at $T=24$: revin-DLinear (left) and the
learnable-gate variant (right). The gate variant's error panel is visibly
wider, consistent with Table~\ref{tab:main}.}
\label{fig:baselines}
\end{figure}

\section{Verified insertion points for follow-up research}
\label{sec:checkpoints}

A key purpose of this reproduction is to certify that the framework can
host new experiments with minimal friction. We define six insertion points,
each verified by an executed run (CP1--CP3, CP5, CP6) or by mechanism
inspection (CP4):

\begin{enumerate}
  \item \textbf{CP1 --- New model.} The registry auto-discovers any
        \texttt{models/*.py} defining \texttt{class Model(nn.Module)} taking
        \texttt{configs}. A new architecture is one file plus one
        \texttt{run.py} invocation.
        \texttt{--model <name>}
  \item \textbf{CP2 --- New data.} Any CSV with columns
        \texttt{date,Open,High,Low,Close} in \texttt{dataset/};
        adapt \texttt{--data\_path} and \texttt{--enc\_in}. Verified with
        freshly fetched GSPC data.
  \item \textbf{CP3 --- Fusion operator.} The gate variants
        (\texttt{-gatev1/-gatev2/-gatev3}) are drop-in comparators under an
        identical protocol; the controlled-comparison design transfers
        directly to any new fusion/routing operator.
        \texttt{--model DeReFusion-gatev2-learnable}
  \item \textbf{CP4 --- Ablations.} \texttt{-woDy}, \texttt{-woLSTM},
        \texttt{-woTransformer} variants isolate component contributions
        without protocol changes.
  \item \textbf{CP5 --- Sweep harness.} Horizon grid ($T\in\{1,7,12,24,36\}$),
        look-back ($L\in\{48,96,192\}$), and capacity ($d_{\mathrm{model}}$)
        are single flags; the batch runner supports resumable sweeps.
  \item \textbf{CP6 --- Metrics and figures.} Every run emits six metrics,
        parameter counts, training/inference timings, and six
        publication-style figures from saved arrays.
\end{enumerate}

For structure-aware expert routing, the minimal experiment is therefore:
implement the router replacement as CP1, compare against CP3's fusion
family under CP5's sweep, and report via CP6 --- reusing the paper's
fairness argument wholesale.

\section{Engineering lessons}
\label{sec:lessons}

The following pitfalls were encountered and fixed; they are recorded so
that future runs are unblocked from the start.

\begin{enumerate}
  \item \textbf{Hidden dependency chain.} Beyond the README's requirements,
        \texttt{patool}, \texttt{huggingface\_hub}, \texttt{sktime}
        (with \texttt{scikit-base}), and \texttt{datasets} are imported
        unconditionally by the data pipeline; each missing module aborts
        startup.
  \item \textbf{Console encoding.} An emoji in the lazy-loading log line
        crashes Windows GBK consoles (\texttt{UnicodeEncodeError}); fixed by
        replacing it with an ASCII marker (committed to the fork).
        Setting \texttt{PYTHONIOENCODING=utf-8} is a robust belt-and-braces
        measure.
  \item \textbf{CPU mode.} \texttt{run.py} defaults to CUDA and asserts when
        torch lacks CUDA support; \texttt{--no\_use\_gpu} is required on
        CPU-only machines.
  \item \textbf{Dependency pins.} \texttt{sktime} requires
        \texttt{joblib<1.6}; when the resolver fights an installed newer
        \texttt{joblib}, pin \texttt{joblib==1.5.3} explicitly. Wheel
        downloads that stall can be recovered with
        \texttt{pip download} (resumable) plus local installation.
  \item \textbf{Dataset in version control.} The upstream repository
        ignores \texttt{dataset/}; the data must be refetched. Our fork
        force-adds the ten CSVs so the repository is self-contained.
  \item \textbf{Non-determinism.} \texttt{run.py} seeds Python/NumPy/PyTorch
        but not cuDNN; with CPU training and a single seed our runs are
        deterministic, but multi-seed reporting remains necessary on CUDA.
\end{enumerate}

\section{Conclusion}
\label{sec:conclusion}

We reproduced the DeReFusion pipeline end to end from the official code with
newly fetched data and confirmed its central qualitative conclusion:
under a controlled protocol, parameter-free additive fusion outperforms
learnable gating, and a ${\sim}28$k-parameter hybrid model is competitive
with the strongest linear baseline. The reproduction doubles as a verified
experiment harness: six documented insertion points allow follow-up
research --- structure-aware expert routing for financial panels --- to be
evaluated under the same fairness protocol at minimal engineering cost.

\section*{CRediT authorship contribution statement}
\textbf{Dongxu Jiang}: Writing -- original draft, Investigation, Formal
analysis, Software. (Supervisor and collaborator contributions to be
finalized at submission.)

\section*{Declaration of competing interest}
The authors declare that they have no known competing financial interests
or personal relationships that could have appeared to influence the work
reported in this paper.

\section*{Data availability}
All data used are publicly available from Yahoo Finance. The fetch script,
forked code, dataset CSVs, and per-experiment artifacts are available at
the authors' public repository fork.

\bibliographystyle{elsarticle-num}
\begin{thebibliography}{10}
\bibitem{hsieh2026derefusion}
C.-C. Hsieh, M.-Y. Chen,
DeReFusion: A controlled comparison of soft computing fusion strategies for
financial time series forecasting via a decomposition-residual architecture,
Applied Soft Computing 203 (2026) 116252.
\href{https://doi.org/10.1016/j.asoc.2026.116252}{doi:10.1016/j.asoc.2026.116252}.

\bibitem{kim2022revin}
T. Kim, J. Kim, Y. Tae, C. Park, J. Choi, G. Choo,
Reversible instance normalization for accurate time-series forecasting
against distribution shift,
in: International Conference on Learning Representations (ICLR), 2022.

\bibitem{zeng2023dlinear}
A. Zeng, M. Chen, L. Zhang, Q. Xu,
Are transformers effective for time series forecasting?,
in: AAAI Conference on Artificial Intelligence, 2023.

\bibitem{vaswani2017attention}
A. Vaswani, N. Shazeer, N. Parmar, J. Uszkoreit, L. Jones, A.N. Gomez,
{\L}. Kaiser, I. Polosukhin,
Attention is all you need,
in: Advances in Neural Information Processing Systems (NeurIPS), 2017.

\bibitem{hochreiter1997lstm}
S. Hochreiter, J. Schmidhuber,
Long short-term memory,
Neural Computation 9 (1997) 1735--1780.

\bibitem{tslib}
H. Wu, T. Hu, Y. Liu, H. Zhou, J. Wang, M. Long,
TimesNet: Temporal 2D-variation modeling for general time series analysis,
in: International Conference on Learning Representations (ICLR), 2023.
(Time-Series-Library code base.)

\end{thebibliography}

\end{document}
